\section{Assumptions and Limitations}
\label{sec:assumptions}

The models in \cref{sec:simulation-model,sec:cornering-model,sec:track-geometry,sec:cda-fitting}
make a number of simplifying assumptions. They are listed here plainly, so that
their limitations are clear to anyone interpreting the simulation's output or a
fitted \CdA{} value.

\begin{itemize}
    \item \textbf{Point-mass dynamics; quasi-static lean, no traction check.}
    The rider and bicycle are treated as a single point mass moving along the
    track's arc length. Lean angle is an algebraic function of instantaneous
    speed and curvature (\cref{sec:quasi-static-lean}), not a dynamical state
    with its own inertia, and is computed unconditionally from speed and
    curvature with no friction-circle/traction-limit check on whether the tire
    could actually supply the implied lateral friction.

    \item \textbf{Idealized (as-if-banked) lean geometry.} The
    rolling-resistance normal-force increase of \cref{eq:normal-force} treats
    leaning as equivalent to riding a surface banked at angle $\theta$ (no
    lateral tire friction needed to supply the centripetal force), rather than
    modeling the actual friction-based lean mechanics of a bicycle on a flat
    track (\cref{sec:rolling-resistance}).

    \item \textbf{CdA fit approximates wheel speed with center-of-mass speed.}
    The \CdA{} fit's distance-from-speed integration and lean-rate chain rule
    (\cref{sec:cda-dvdt}) use the recorded (center-of-mass) speed directly,
    rather than the wheel-path speed $v_{\mathrm{wheel}}$ of
    \cref{eq:wheel-speed} that the forward simulation uses for the same
    purpose; the difference is second-order in $\kappa h \sin\theta$ and small
    for realistic lean angles, but is not corrected for.

    \item \textbf{Flat-track assumption.} Neither the forward simulation nor the
    \CdA{} fit includes a gravitational grade term. This is appropriate for a
    velodrome or flat closed-circuit use case, but a limitation for road data
    with meaningful elevation change.

    \item \textbf{No wind-speed correction.} \texttt{CdAFitter} uses ground
    speed directly in the drag equation (\cref{eq:aero-drag}); wind speed is
    present in recorded \texttt{.fit} telemetry but is not subtracted from or
    added to the speed used in the fit.

    \item \textbf{Steady environmental and mechanical conditions.} Rolling
    resistance coefficient, air density, and drivetrain mechanical loss are each
    taken as a single constant value across an entire simulation run or fitting
    window, not varied over time or estimated from the data.

    \item \textbf{Unsmoothed finite-difference acceleration in the CdA fit.} The
    acceleration of \cref{sec:cda-dvdt} is a raw central difference of the
    recorded speed trace, with no low-pass filtering beforehand; noisy speed
    samples propagate directly into a noisy \CdA{} estimate.

    \item \textbf{Unweighted least squares.} The regression of
    \cref{eq:cda-closed-form} weights every sample equally and includes no
    outlier-robust reweighting, so a small number of corrupted samples can shift
    the fit disproportionately.

    \item \textbf{Idealized track geometry.} The track geometry of
    \cref{sec:track-geometry} is a two-straight, two-semicircle oval with a
    box-filter-smoothed (not clothoid) transition; it approximates, but does not
    exactly reproduce, the geometry of a real track or velodrome.

    \item \textbf{Position reconstruction is visualization-only.} The
    Cartesian centerline computed in \cref{sec:position-reconstruction} is used
    only to draw the track shape in the application; it has no influence on the
    physics of \cref{sec:simulation-model,sec:cornering-model}.
\end{itemize}
